\section{Group algebra, class functions, and symmetric algebras}
\begin{definition}[Group algebra]
    Let \(G\) be a finite group, and let \(R\) be a commutative unital ring.
    The group algebra \(RG\) is the free \(R\)-module with basis \(G\), whose
    multiplication is determined by the group multiplication of \(G\).

    Thus every element of \(RG\) can be written uniquely in the form
    \[
        x=\sum_{s\in G} x_s s,
        \qquad
        y=\sum_{t\in G} y_t t,
    \]
    where \(x_s,y_t\in R\). Their product is given by
    \[
        xy
        =
        \sum_{s,t\in G} x_s y_t st
        =
        \sum_{g\in G}
        \left(
        \sum_{\substack{s,t\in G\\ st=g}} x_s y_t
        \right) g.
    \]
\end{definition}

From now on, let \(k\) be a field, and consider the group algebra \(kG\).

\begin{definition}
    Define a \(k\)-linear form
    \[
        \pi:kG\longrightarrow k
    \]
    by
    \[
        \pi(g)=
        \begin{cases}
            1, & g=1,     \\
            0, & g\neq 1,
        \end{cases}
    \]
    for all \(g\in G\). Equivalently, \(\pi(x)\) is the coefficient of the identity
    element \(1\in G\) in \(x\in kG\).
\end{definition}

\begin{lemma}
    Let \(x,y\in kG\). Then
    \[
        \pi(xy)=\pi(yx).
    \]
\end{lemma}

\begin{proof}
    Write
    \[
        x=\sum_{s\in G}x_s s,
        \qquad
        y=\sum_{t\in G}y_t t.
    \]
    Then the coefficient of \(1\in G\) in \(xy\) is
    \[
        \pi(xy)
        =
        \sum_{\substack{s,t\in G\\ st=1}}x_s y_t
        =
        \sum_{s\in G}x_s y_{s^{-1}}.
    \]
    Similarly,
    \[
        \pi(yx)
        =
        \sum_{\substack{t,s\in G\\ ts=1}}y_t x_s
        =
        \sum_{t\in G}y_t x_{t^{-1}}.
    \]
    After replacing \(t\) by \(s^{-1}\), we get
    \[
        \pi(yx)
        =
        \sum_{s\in G}y_{s^{-1}}x_s
        =
        \sum_{s\in G}x_s y_{s^{-1}}
        =
        \pi(xy),
    \]
    because \(k\) is commutative.
\end{proof}

\begin{lemma}
    The map
    \[
        \psi:kG\longrightarrow \operatorname{Hom}_k(kG,k),
        \qquad
        x\longmapsto \bigl(y\mapsto \pi(xy)\bigr),
    \]
    is a \(k\)-linear isomorphism.
\end{lemma}

\begin{proof}
    It is clear that \(\psi\) is \(k\)-linear.

    Let \(f\in \operatorname{Hom}_k(kG,k)\). We claim that the inverse image of
    \(f\) under \(\psi\) is
    \[
        \psi^{-1}(f)
        =
        \sum_{g\in G} f(g^{-1})g.
    \]
    Indeed, set
    \[
        x=\sum_{g\in G} f(g^{-1})g.
    \]
    For any \(h\in G\), we have
    \[
        \psi(x)(h)
        =
        \pi(xh)
        =
        \pi\left(
        \sum_{g\in G} f(g^{-1})gh
        \right).
    \]
    The coefficient of \(1\) occurs precisely when \(gh=1\), that is,
    \(g=h^{-1}\). Therefore
    \[
        \psi(x)(h)
        =
        f\bigl((h^{-1})^{-1}\bigr)
        =
        f(h).
    \]
    Since \(G\) is a basis of \(kG\), this proves \(\psi(x)=f\). Hence \(\psi\) is
    surjective. The same formula also shows uniqueness of \(x\), so \(\psi\) is
    injective. Therefore \(\psi\) is a \(k\)-linear isomorphism.
\end{proof}

\begin{definition}[Class function]
    A function
    \[
        f:G\longrightarrow k
    \]
    is called a class function if
    \[
        f(gxg^{-1})=f(x)
    \]
    for all \(g,x\in G\).
\end{definition}

\begin{remark}
    If \(f:G\to k\) is a class function, then
    \[
        f(xy)=f(yx)
    \]
    for all \(x,y\in G\). Indeed,
    \[
        xy=x(yx)x^{-1},
    \]
    so \(xy\) and \(yx\) are conjugate.

    Extending \(f\) linearly, we obtain a \(k\)-linear map
    \[
        f:kG\longrightarrow k.
    \]
    Thus class functions on \(G\) may be regarded as special linear forms on
    \(kG\).
\end{remark}

\begin{remark}
    There is a natural identification
    \[
        \{\,\text{functions }G\to k\,\}
        \cong
        \operatorname{Hom}_k(kG,k)
        =
        (kG)^*.
    \]
    Indeed, since \(G\) is a \(k\)-basis of \(kG\), a linear form on \(kG\) is
    uniquely determined by its values on the basis elements \(g\in G\).
\end{remark}

\begin{definition}[Symmetric algebra]
    Let \(A\) be a finite-dimensional \(k\)-algebra. We say that \(A\) is a
    \emph{symmetric algebra} if there exists a \(k\)-bilinear form
    \[
        \langle -,-\rangle:A\times A\longrightarrow k
    \]
    such that:

    \begin{enumerate}
        \item \(\langle -,-\rangle\) is nondegenerate;
        \item \(\langle a,b\rangle=\langle b,a\rangle\) for all \(a,b\in A\);
        \item \(\langle ab,c\rangle=\langle a,bc\rangle\) for all \(a,b,c\in A\).
    \end{enumerate}
\end{definition}

\begin{definition}[Equivalent definition]
    A finite-dimensional \(k\)-algebra \(A\) is called symmetric if
    \[
        A\cong A^*
    \]
    as \(A\)-bimodules, where
    \[
        A^*=\operatorname{Hom}_k(A,k).
    \]
\end{definition}

\begin{proposition}[Equivalent definitions of symmetric algebras]
    Let \(A\) be a finite-dimensional unital associative \(k\)-algebra. Put
    \[
        A^*=\operatorname{Hom}_k(A,k),
    \]
    and give \(A^*\) the \(A\)-bimodule structure
    \[
        (a\cdot f\cdot b)(x)
        =
        f(bxa),
        \qquad a,b,x\in A,\ f\in A^*.
    \]
    Then the following are equivalent:

    \begin{enumerate}
        \item[\textup{(i)}] There exists a nondegenerate \(k\)-bilinear form
              \[
                  \langle -,-\rangle:A\times A\longrightarrow k
              \]
              such that
              \[
                  \langle a,b\rangle=\langle b,a\rangle
              \]
              for all \(a,b\in A\), and
              \[
                  \langle ab,c\rangle=\langle a,bc\rangle
              \]
              for all \(a,b,c\in A\).

        \item[\textup{(ii)}] There is an isomorphism
              \[
                  A\cong A^*
              \]
              as \(A\)-bimodules.
    \end{enumerate}
\end{proposition}

\begin{proof}
    Assume first that \(\langle -,-\rangle\) is a nondegenerate symmetric
    associative bilinear form on \(A\). Define
    \[
        \Phi:A\longrightarrow A^*,
        \qquad
        \Phi(a)(x)=\langle x,a\rangle .
    \]
    Since the bilinear form is nondegenerate and \(A\) is finite-dimensional,
    \(\Phi\) is a \(k\)-linear isomorphism.

    We claim that \(\Phi\) is an \(A\)-bimodule homomorphism. Let
    \(a,c,d,x\in A\). Then
    \[
        \begin{aligned}
            (c\cdot \Phi(a)\cdot d)(x)
             & =\Phi(a)(dxc)         \\
             & =\langle dxc,a\rangle \\
             & =\langle dx,ca\rangle \\
             & =\langle d,xca\rangle \\
             & =\langle xca,d\rangle \\
             & =\langle xc,ad\rangle \\
             & =\langle x,cad\rangle \\
             & =\Phi(cad)(x).
        \end{aligned}
    \]
    Therefore
    \[
        \Phi(cad)=c\cdot \Phi(a)\cdot d.
    \]
    Hence \(\Phi:A\to A^*\) is an isomorphism of \(A\)-bimodules.

    Conversely, assume that there is an \(A\)-bimodule isomorphism
    \[
        \Phi:A\longrightarrow A^*.
    \]
    Define a \(k\)-bilinear form on \(A\) by
    \[
        \langle x,y\rangle=\Phi(y)(x).
    \]
    Since \(\Phi\) is a \(k\)-linear isomorphism, this bilinear form is
    nondegenerate.

    We first prove associativity. For \(a,b,c\in A\), since \(\Phi\) is a left
    \(A\)-module homomorphism, we have
    \[
        \Phi(bc)=b\cdot \Phi(c).
    \]
    Thus
    \[
        \begin{aligned}
            \langle a,bc\rangle
             & =\Phi(bc)(a)           \\
             & =(b\cdot \Phi(c))(a)   \\
             & =\Phi(c)(ab)           \\
             & =\langle ab,c\rangle .
        \end{aligned}
    \]
    Hence
    \[
        \langle ab,c\rangle=\langle a,bc\rangle .
    \]

    It remains to prove symmetry. Let
    \[
        \lambda=\Phi(1)\in A^*.
    \]
    Since \(\Phi\) is an \(A\)-bimodule homomorphism, for every \(y\in A\) we have
    \[
        \Phi(y)=\Phi(y\cdot 1)=y\cdot \Phi(1)=y\cdot \lambda
    \]
    and also
    \[
        \Phi(y)=\Phi(1\cdot y)=\Phi(1)\cdot y=\lambda\cdot y.
    \]
    Therefore, for all \(x,y\in A\),
    \[
        \Phi(y)(x)
        =
        (y\cdot \lambda)(x)
        =
        \lambda(xy),
    \]
    and also
    \[
        \Phi(y)(x)
        =
        (\lambda\cdot y)(x)
        =
        \lambda(yx).
    \]
    Thus
    \[
        \lambda(xy)=\lambda(yx)
    \]
    for all \(x,y\in A\). Hence
    \[
        \begin{aligned}
            \langle x,y\rangle
             & =\Phi(y)(x)           \\
             & =\lambda(xy)          \\
             & =\lambda(yx)          \\
             & =\Phi(x)(y)           \\
             & =\langle y,x\rangle .
        \end{aligned}
    \]
    So the bilinear form is symmetric.

    Therefore the existence of a nondegenerate symmetric associative bilinear form
    on \(A\) is equivalent to the existence of an \(A\)-bimodule isomorphism
    \[
        A\cong A^*.
    \]
\end{proof}

\begin{corollary}
    The group algebra \(kG\) is a symmetric algebra.
\end{corollary}

\begin{proof}
    The bilinear form
    \[
        \langle x,y\rangle
        =
        \pi(xy)
    \]
    is symmetric by the first lemma:
    \[
        \langle x,y\rangle
        =
        \pi(xy)
        =
        \pi(yx)
        =
        \langle y,x\rangle.
    \]
    The second lemma says precisely that the induced map
    \[
        kG\longrightarrow (kG)^*,
        \qquad
        x\longmapsto \langle x,-\rangle,
    \]
    is a \(k\)-linear isomorphism. Therefore the bilinear form is nondegenerate,
    and hence \(kG\) is a symmetric algebra.
\end{proof}

Let
\[
    F(G,k)
\]
denote the \(k\)-vector space of all functions \(G\to k\). We identify
\(F(G,k)\) with \((kG)^*=\operatorname{Hom}_k(kG,k)\) by extending functions
linearly from \(G\) to \(kG\).

Let
\[
    \operatorname{CF}(G,k)
\]
denote the subspace of class functions on \(G\), that is,
\[
    \operatorname{CF}(G,k)
    =
    \{\,f:G\to k \mid f(hgh^{-1})=f(g)
    \text{ for all } g,h\in G\,\}.
\]

Recall that for \(f\in F(G,k)\), we define
\[
    f^\circ
    =
    \sum_{g\in G} f(g^{-1})g
    \in kG.
\]

\begin{lemma}
    Let \(a\in kG\) and \(f\in F(G,k)\). Define \(a\cdot f\in F(G,k)\) by
    \[
        (a\cdot f)(x)=f(xa),
        \qquad x\in kG.
    \]
    Then
    \[
        (a\cdot f)^\circ=a f^\circ .
    \]
\end{lemma}

\begin{proof}
    Write
    \[
        a=\sum_{h\in G}a_hh.
    \]
    Then
    \[
        \begin{aligned}
            (a\cdot f)^\circ
             & =\sum_{g\in G}(a\cdot f)(g^{-1})g            \\
             & =\sum_{g\in G}f(g^{-1}a)g                    \\
             & =\sum_{g\in G}\sum_{h\in G}a_h f(g^{-1}h)g .
        \end{aligned}
    \]
    On the other hand,
    \[
        \begin{aligned}
            af^\circ
             & =
            \left(\sum_{h\in G}a_hh\right)
            \left(\sum_{t\in G}f(t^{-1})t\right) \\
             & =
            \sum_{h,t\in G}a_h f(t^{-1})ht.
        \end{aligned}
    \]
    Putting \(g=ht\), equivalently \(t=h^{-1}g\), we get
    \[
        t^{-1}=g^{-1}h.
    \]
    Hence
    \[
        af^\circ
        =
        \sum_{g\in G}\sum_{h\in G}a_h f(g^{-1}h)g
        =
        (a\cdot f)^\circ .
    \]
\end{proof}

Let
\[
    \operatorname{Cl}(G)
\]
denote the set of conjugacy classes of \(G\). For each
\(C\in \operatorname{Cl}(G)\), define the characteristic function
\[
    \gamma_C:G\longrightarrow k
\]
by
\[
    \gamma_C(g)
    =
    \begin{cases}
        1, & g\in C,    \\
        0, & g\notin C.
    \end{cases}
\]
Also define the class sum
\[
    \widehat{C}
    =
    \sum_{g\in C}g
    \in kG.
\]

For a class function \(f\in \operatorname{CF}(G,k)\), if \(g_C\in C\) is a
chosen representative for each conjugacy class \(C\), then
\[
    f
    =
    \sum_{C\in \operatorname{Cl}(G)} f(g_C)\gamma_C.
\]
Moreover,
\[
    f^\circ
    =
    \sum_{C\in \operatorname{Cl}(G)} f(g_C^{-1})\widehat{C}.
\]

\begin{lemma}
    Let \(Z(kG)\) denote the center of the group algebra \(kG\), that is,
    \[
        Z(kG)
        =
        \{\,x\in kG\mid xy=yx \text{ for all } y\in kG\,\}.
    \]
    Then:

    \begin{enumerate}
        \item The family
              \[
                  \{\,\gamma_C\mid C\in \operatorname{Cl}(G)\,\}
              \]
              is a basis of \(\operatorname{CF}(G,k)\).

        \item The family
              \[
                  \{\,\widehat{C}\mid C\in \operatorname{Cl}(G)\,\}
              \]
              is a basis of \(Z(kG)\).

        \item Hence
              \[
                  \dim_k \operatorname{CF}(G,k)
                  =
                  |\operatorname{Cl}(G)|.
              \]

        \item The map
              \[
                  F(G,k)\longrightarrow kG,
                  \qquad
                  f\longmapsto f^\circ,
              \]
              restricts to a \(k\)-linear isomorphism
              \[
                  \operatorname{CF}(G,k)
                  \xrightarrow{\sim}
                  Z(kG).
              \]
              In particular,
              \[
                  \dim_k Z(kG)=|\operatorname{Cl}(G)|.
              \]
    \end{enumerate}
\end{lemma}

\begin{proof}
    First, every class function is constant on each conjugacy class. Therefore, if
    \(g_C\in C\) is a chosen representative for each \(C\in \operatorname{Cl}(G)\),
    then every \(f\in \operatorname{CF}(G,k)\) can be written uniquely as
    \[
        f
        =
        \sum_{C\in \operatorname{Cl}(G)} f(g_C)\gamma_C.
    \]
    Thus
    \[
        \{\,\gamma_C\mid C\in \operatorname{Cl}(G)\,\}
    \]
    is a basis of \(\operatorname{CF}(G,k)\). In particular,
    \[
        \dim_k\operatorname{CF}(G,k)=|\operatorname{Cl}(G)|.
    \]

    Next, we prove that the class sums form a basis of the center. Let
    \[
        x=\sum_{g\in G}x_g g\in kG.
    \]
    Then \(x\in Z(kG)\) if and only if
    \[
        hxh^{-1}=x
    \]
    for every \(h\in G\). But
    \[
        hxh^{-1}
        =
        \sum_{g\in G}x_g hgh^{-1}.
    \]
    Therefore \(hxh^{-1}=x\) for all \(h\in G\) if and only if the coefficient
    \(x_g\) is constant on conjugacy classes. Hence every central element has a
    unique expression of the form
    \[
        x
        =
        \sum_{C\in \operatorname{Cl}(G)} a_C \widehat{C},
        \qquad a_C\in k.
    \]
    Thus
    \[
        \{\,\widehat{C}\mid C\in \operatorname{Cl}(G)\,\}
    \]
    is a basis of \(Z(kG)\).

    Finally, for \(C\in \operatorname{Cl}(G)\), we have
    \[
        \gamma_C^\circ
        =
        \sum_{g\in G}\gamma_C(g^{-1})g
        =
        \sum_{g^{-1}\in C}g
        =
        \sum_{g\in C^{-1}}g
        =
        \widehat{C^{-1}}.
    \]
    Since the map
    \[
        C\longmapsto C^{-1}
    \]
    is a permutation of the set of conjugacy classes, the map
    \[
        f\longmapsto f^\circ
    \]
    sends the basis
    \[
        \{\,\gamma_C\mid C\in \operatorname{Cl}(G)\,\}
    \]
    of \(\operatorname{CF}(G,k)\) onto the basis
    \[
        \{\,\widehat{C}\mid C\in \operatorname{Cl}(G)\,\}
    \]
    of \(Z(kG)\). Therefore it restricts to a \(k\)-linear isomorphism
    \[
        \operatorname{CF}(G,k)
        \xrightarrow{\sim}
        Z(kG).
    \]
    Consequently,
    \[
        \dim_k Z(kG)=|\operatorname{Cl}(G)|.
    \]
\end{proof}
